% Euler's Equations: Symbolic vs Numerical Mathematics
% Claes Johnson, sole author. DRAFT -- summing up.
\documentclass[11pt]{article}
\usepackage[margin=1.15in]{geometry}
\usepackage{amsmath,amssymb}
\usepackage{url}
\usepackage[hidelinks]{hyperref}
\setlength{\parskip}{3pt}

\newcommand{\grad}{\nabla}
\newcommand{\eps}{\varepsilon}

\title{\vspace{-2em}Euler's Equations:\\Symbolic versus Numerical Mathematics}
\author{Claes Johnson\\[0.4em]
\small KTH Royal Institute of Technology, Stockholm\\
\small \texttt{claesjohnson@gmail.com}}
\date{\today}

\begin{document}
\maketitle

\begin{abstract}
\noindent
Three things are to be kept apart: a fluid flow, which is a phenomenon; Euler's equations, which are
a model of it; and a proof about the equations, which is a statement about the model. We state the
equations, note that they are parameter free, and observe that the mathematical question asked of
them --- existence, regularity, uniqueness --- has not been answered by analysis in symbolic form,
because the solutions turn out to have a complicated nature which symbolic analysis has no grip on.
Solutions in digital form over a mesh can be computed by a model we name \emph{Numerical Euler} ---
a stabilised finite element discretisation together with duality-based output error control and
adaptive refinement --- and they are turbulent, fluctuating over a range of scales, with well
determined mean values such as drag and lift which agree with measurement and arrive with a computed
bound on their error. We
then ask what a recent symbolic proof of finite-time blowup establishes, and conclude that it is a
statement about the model which does not reach either of the other two levels: it exhibits no
individual case, so nothing can be inserted into the equations, inspected, or measured.
\end{abstract}

\section{Euler, 1755}\label{sec:eq}

In 1755 Leonhard Euler presented to the Berlin Academy the memoir \emph{Principes g\'en\'eraux du
mouvement des fluides}~\cite{euler1755}, in which he wrote down the equations of motion of an
incompressible fluid. He took them to contain the whole of the mechanics of fluids, and he was
explicit that what then stood in the way of using them was not physics but mathematics:

\begin{quote}\small
And if it is not permitted to us to penetrate to a complete knowledge concerning the motion of
fluids, it is not to mechanics, or to the insufficiency of the known principles of motion, that we
must attribute the cause. It is analysis itself which abandons us here.
\end{quote}

\noindent What Euler had in mind was \emph{ideal} fluid flow --- inviscid incompressible flow --- as
a model of the slightly viscous, nearly incompressible flow of real fluids such as water and air at
subsonic speeds. He formulated his model in the following form, named \emph{Euler's equations}:
\begin{align}
\dot u + (u\cdot\grad)u + \grad p &= f \qquad\text{in }\Omega, \label{eq:mom}\\[2pt]
\grad\cdot u &= 0 \qquad\;\;\text{in }\Omega, \label{eq:inc}\\[2pt]
u\cdot n &= 0 \qquad\;\;\text{on }\Gamma . \label{eq:slip}
\end{align}
Here $\Omega\subset\mathbb{R}^3$ is the domain occupied by the fluid and $\Gamma$ its boundary,
$x=(x_1,x_2,x_3)$ is a Euclidean coordinate and $t$ time; $u(x,t)\in\mathbb{R}^3$ is the velocity and
$p(x,t)$ the pressure of a fluid of unit density; $f$ is a given force, $n$ the outward normal to
$\Gamma$, and $u(\cdot,0)=u^0$ a given initial velocity. The equations hold for $0<t\le T$.

What they say is the following. Equation~\eqref{eq:mom} is Newton's second law: for unit density the
mass of a fluid particle is its volume, so force equals acceleration directly, the left-hand side
being the acceleration of a particle following the flow and the right-hand side the force on it ---
the pressure gradient, which is what the surrounding fluid exerts, together with any body force.
Equation~\eqref{eq:inc} states incompressibility, the volume of a particle not changing, with the
pressure as whatever field enforces it. Condition~\eqref{eq:slip} states that the fluid does not pass
through the wall, and nothing more: the normal component is fixed by impermeability, nothing in the
equations determines the tangential velocity there, and prescribing it as well would over-determine
the problem. There is no viscosity anywhere, and therefore no boundary layer.

\paragraph{The Eulerian representation is a mathematical invention.} The equations look simple and
innocent, and they are neither. They embody a fundamental innovation of mathematical rather than
physical nature, in what is called the \emph{Eulerian} representation. A fictional observer is posted
at each point $(x,t)$ of space-time, and records the velocity $u(x,t)$ of whatever fluid particle is
passing at that instant, together with the pressure $p(x,t)$ there. No such observer exists in the
fluid. Nothing physical sits at a fixed point in space watching particles go by; the fluid consists
of the particles.

The physical description is the other one. In the \emph{Lagrangian} representation the individual
particles are followed along their paths, each governed directly by Newton's second law, which is
what the mechanics actually says. That description is the faithful one and it is close to
unusable for computation: the particles go where they go, the mesh one would have to carry with them
deforms without bound, and in a turbulent flow it is destroyed within a few time units.

The Eulerian description, which corresponds to nothing physical, is what makes efficient numerics
possible --- a fixed domain, a fixed mesh, fields evaluated at fixed points, and the particle motion
appearing only through the convective term $(u\cdot\grad)u$. This is worth noticing early, because
the paper is about what symbolic mathematics contributes. Here it contributes the representation
itself: a mathematical fiction with no physical counterpart, invented in 1755, without which the
computations of Section~\ref{sec:fenics} could not be carried out at all.

\paragraph{The model is parameter free.} There is no constant in
\eqref{eq:mom}--\eqref{eq:slip} to be measured, fitted or chosen. What enters as data is the
\emph{geometry} --- the domain, which is to say the shape of the body the fluid flows past --- the
\emph{forcing}, and the initial state. Nothing else. Whatever the equations predict is a consequence
of the geometry and the forcing alone, and any disagreement with observation is a fault of the model
and cannot be repaired by adjustment.

\paragraph{Physics without physical input.} One consequence of that deserves to be stated at the
outset, since it is easy to pass over. The drag and the lift of a body of \emph{arbitrary shape} can
be computed from \eqref{eq:mom}--\eqref{eq:slip} alone. No measurement enters, no coefficient is
looked up, and nothing shape-specific is supplied: one states the geometry and solves. A correct
statement about the physical world is thereby obtained with no input from the physical world beyond
Newton's second law and the incompressibility of the fluid --- which were in hand before the body was
drawn.

It is instructive to set that against the hydraulics of Section~\ref{sec:dalembert}, where every
shape requires its own measured coefficient and nothing can be said about a body before it is built
and tested. Here the drag of a shape nobody has ever constructed follows from the equations, and the
place where this is put to the test is the workshops of Section~\ref{sec:fenics}, at which
predictions are submitted \emph{before} the measurements are released.

The character of the knowledge so obtained is worth a word. It is close to an \emph{a priori} truth
reached by a thought experiment, in the spirit in which Einstein used them: the conclusion about the
world is forced by principles already accepted, and the experiment is performed in the equations
rather than in the laboratory. One should not overstate it --- Newton's law and incompressibility are
themselves empirical, and the parameter-free claim is that nothing \emph{further} is needed. But
granted those, the resistance of any shape whatever is a matter of computation and not of
measurement, and that is a stronger thing than fluid mechanics was taken to possess for most of its
history.

\section{d'Alembert's paradox, and the split}\label{sec:dalembert}

The expectation that the mechanics of fluids was contained in two short mathematical equations met an
objection which was in print before the equations were. Three years earlier, in the \emph{Essai d'une
nouvelle th\'eorie de la r\'esistance des fluides} of 1752~\cite{dalembert1752}, d'Alembert had shown
that a body moving steadily through an ideal fluid experiences \emph{no resistance at all}. Ships
meet resistance and cannonballs are slowed, so the conclusion is absurd --- and it is a conclusion
about precisely the solutions of \eqref{eq:mom}--\eqref{eq:slip}, to which it duly attached itself.

It deepened rather than dissolved. In the nineteenth century complex analysis supplied a machinery
generating ideal flow as \emph{potential flow} for any body one cared to draw, all of it smooth, all
of it elegant, and all of it with zero drag. The better the mathematics got, the larger grew the
family of solutions describing nothing. Something was wrong with the equations and nobody could say
what.

\paragraph{What it cost.} The consequence was not a puzzle set aside for later. It was a divorce.
Mathematical fluid mechanics, unable to predict the resistance of a body moving through water, became
a subject whose central prediction was known to be false --- which is to say a joke, and it was
treated as one by the people who had to build ships and, later, aircraft. The practical science went
its own way as \emph{hydraulics}: correlations, loss coefficients, model tests in towing tanks and
wind tunnels, and a working knowledge of turbulent flow that owed nothing to Euler's equations and
could not be derived from them. The division is well known and was put best as a distinction between
engineers who observed what could not be explained and mathematicians who explained what could not be
observed.\footnote{The remark is usually credited to Hinshelwood and is quoted by Lighthill.
[TO VERIFY.]}

So the two halves of the subject parted. Real fluids were understood empirically and without
mathematics; the mathematics was exact, beautiful and about nothing.

\section{Leray, 1934: turbulent solutions without uniqueness}\label{sec:leray}

Nothing in the mathematics moved this position for the better part of two centuries. What moved it
was Leray, in the paper of 1934 on the motion of a viscous liquid filling space~\cite{leray}. Leray
abandoned the requirement that a solution be smooth and defined a \emph{weak} solution, satisfying
the equations in an integrated sense against test functions. For such solutions he proved global
existence. And he named them for what he took them to represent: \emph{solutions turbulentes}.

This was the decisive step and it did not finish the job. What Leray obtained was existence without
\emph{uniqueness}, and uniqueness is what a model needs if it is to predict anything. A theory which
says that solutions exist, and cannot say which one a given body in a given stream will produce,
determines no drag and no lift. The subject gained a solution concept adequate to turbulence and
still could not deliver a number, and the uniqueness question has remained open since --- for weak
solutions of Euler it is now known to be \emph{false}.

\section{Numerical Euler: weak solutions with residual control}\label{sec:digital}

The gap Leray left is where the present account begins. It was closed, in our view, not by proving
uniqueness of the solution --- which is unavailable, and false --- but by changing what uniqueness is
asked of. The solution concept set out in~\cite{book} is a weak solution together with
\emph{control of the residual}, and what it delivers is \emph{uniqueness of mean-value outputs}: the
flow field is not determined \emph{pointwise in space and time}, but the drag and the lift, as
\emph{mean-value outputs}, are well determined quantities.

We name the computational model \textbf{Numerical Euler}. It has three parts and the third is not an
afterthought:

\begin{enumerate}
\item the equations \eqref{eq:mom}--\eqref{eq:slip}, discretised by a stabilised finite element
method over a mesh~\cite{book}, the stabilisation acting as a vanishing artificial viscosity, so that
what is computed is a viscosity solution of Euler, the limit taken as the mesh is refined;
\item \emph{duality-based output error control}: for each output of interest --- a drag, a lift ---
the error is bounded by the residual of the computed solution weighted by the solution of an adjoint
problem, which measures how much the residual matters for that output;
\item \emph{adaptivity}: the mesh is refined where the weighted residual says the error lives, and
not uniformly, until the bound on the output is below what is asked.
\end{enumerate}

The error control is part of the model rather than an assessment of it. What Numerical Euler returns
is not a flow but a \emph{pair}: an output and a bound on how wrong that output may be, computed
together. A model returning the first without the second would be a weaker thing, since nothing could
be concluded from it.

\paragraph{And it is a solution of Euler.} That the limit satisfies
\eqref{eq:mom}--\eqref{eq:slip} weakly is not an identification offered because the numbers look
right. The stabilisation is a weighted least-squares control of the residual; the discrete energy
estimate bounds that residual in a weighted norm; testing against a smooth function sends it to zero
in the sense the weak formulation requires. The limit is a weak solution in Leray's sense, as a
theorem --- and unlike Leray's, it comes with a statement of what is determined.

\paragraph{What it produces, none of it put in.} The computed flow is \emph{turbulent}: fluctuating
in time and over a range of scales which widens as the mesh is refined, with no symmetry and no
steady state. Its \emph{mean values are determined}: drag, lift, and the position of the separation
line converge under refinement, each arriving with its bound. And its \emph{velocity stays bounded
while its gradients do not}, the solution being H\"older continuous with exponent about $1/3$ --- the
exponent below which Onsager observed that energy conservation can fail~\cite{onsager} --- and
dissipating energy at a rate independent of the vanishing viscosity.

\paragraph{The viscosity introduced, and why its size does not matter.} A computation at finite
resolution cannot have zero dissipation, so Numerical Euler necessarily introduces a viscosity. This
is the obvious objection to it, and the answer is a computational fact rather than an argument:
\emph{the mean-value outputs do not depend on the absolute size of the viscosity introduced}. Refine
the mesh and the artificial viscosity falls, the resolved scales multiply, the flow field changes in
every detail --- and the drag does not move. The dissipation rate does not fall with the viscosity
either; it settles at a value fixed by the geometry and the free stream.

That last is worth stating as a result and not only as a reassurance about the method. A dissipation
rate which does not vanish with the viscosity is precisely \emph{Kolmogorov's conjecture of finite
dissipation}~\cite{kolmogorov41} --- the zeroth law of turbulence, assumed in the theory of 1941,
observed in experiment since Taylor~\cite{taylor}, and still unproved. The computation supports it
directly: the rate is measured, over a sequence of refinements in which the artificial viscosity
falls by orders of magnitude, and it does not follow the viscosity down. That independence is the
dissipation anomaly in computational form, and it places the computation and the fluid in the same
position rather than in different ones: air has a small viscosity whose absolute size does not
determine the drag of a wing, and Numerical Euler has a small viscosity whose absolute size does not
determine the computed drag. Neither is inviscid and neither needs to be.

\paragraph{d'Alembert's paradox, resolved.} With turbulent solutions available for inspection the
paradox can be answered, and was~\cite{dalembert}. Potential flow is a solution and it is
\emph{unstable} --- stable where the flow attaches on the forward surface and accelerates, unstable
where it separates on the rear surface and is retarded. The instability generates streamwise
vorticity which stretches and breaks down into a turbulent wake, and the wake carries the drag. The
same mechanism applied to a wing gives lift with no circulation condition imposed by
hand~\cite{flight}. The equations were never wrong. What was wrong was the assumption that the
solution of interest is a smooth one, and two centuries of symbolic mathematics had been spent
elaborating exactly the solutions that do not occur.

\section{FEniCS, automation, and the workshops}\label{sec:fenics}

A solution concept is of no use to an engineer who cannot compute with it, and the step from the
concept to a tool was automation. The FEniCS project~\cite{fenics} automates the step from a
variational form to compiled code: the equations are stated, the discretisation is generated, and the
adjoint problem required by the error control is generated with them. What had been a research
computation became a method that could be run by somebody who had not written it.

The result has been taken into use, and the place to see whether it works is not our own publications
but the benchmark workshops, where the cases, the geometries and the measurements belong to the
organisers and predictions are submitted before the measurements are released. Numerical Euler has
been put through them: the high-order CFD workshop case CS1~\cite{cs1}, and validation against the
NASA High Lift Prediction Workshop~4~\cite{ifasd}. Drag within a few per cent, the separation line
where experiment puts it, the lift curve and the stall, the drag crisis~\cite{dragcrisis}.

\paragraph{Three solution concepts.} The arc is a history of what a solution of
\eqref{eq:mom}--\eqref{eq:slip} was taken to be, and it has three stages.

\begin{center}\small
\begin{tabular}{@{}llp{0.40\textwidth}@{}}
\hline
\textbf{potential flow} & pointwise solution & exists, and is \emph{unstable}: determines nothing,
since it does not occur \\[3pt]
\textbf{Leray, 1934} & weak solution & exists globally, \emph{without uniqueness}: determines
nothing, since it does not say which \\[3pt]
\textbf{Numerical Euler} & weak solution with residual control & not determined pointwise in space
and time, but the \emph{mean-value outputs} --- drag, lift --- are well determined, each with a
bound \\
\hline
\end{tabular}
\end{center}

\noindent Each stage weakens the demand made of the solution and strengthens what can be concluded
from it, which is the opposite of the order one expects. The pointwise solution is the strongest
object and settles nothing. The weak solution is a weaker object and settles existence. The weak
solution with its residual controlled is weaker still --- it is not even a single flow --- and it
settles the only question an engineer asks.

That is the end of the arc begun in 1752. The subject which split into a mathematics that explained
what could not be observed and a hydraulics that observed what could not be explained has a model
with no parameters, solved from the equations, delivering the numbers an aircraft designer asks for,
to a tolerance the computation itself states, against measurements held by somebody else.

\paragraph{Symbolic and numerical mathematics together.} None of this is an argument against symbolic
mathematics, and a paper whose first act is to write down three equations is in no position to make
one. The relation is not a rivalry, and it is worth setting out what each contributes, because the
account above has symbolic mathematics doing the decisive work at three separate points.

\emph{It formulates the equations.} Newton's second law is a statement about particles. Writing it
for a continuum, identifying the force of the surrounding fluid as a pressure gradient, and adjoining
incompressibility with the pressure as its multiplier, is symbolic work of the first order, and every
step of it was a judgement about what a fluid is. Three lines with no parameters, from which
everything else here follows.

\emph{It exhibits the nature of the instability.} That potential flow is unstable at separation is
not read off a picture. It is shown by analysis of the \emph{linearised} equations --- a symbolic
calculation, of exactly the classical kind --- which identifies the growing mode, its orientation and
its rate. The computation then confirms what the analysis identifies and shows what becomes of it.
Neither half would do on its own: the linear analysis says what begins and cannot say where it ends,
and the computation shows the end without naming the mechanism.

\emph{It sets up the computation, and interprets it.} The discretisation is derived from the
equations; the residual is defined symbolically; the adjoint problem and the error representation are
derived, not discovered; and the reading of what comes out --- that the gradients scale as they do,
that the dissipation is anomalous in Onsager's sense, that the wake carries the drag --- is
conceptual analysis of computed solutions, which is symbolic mathematics applied to numerical output.

So the two reinforce one another at every stage, and the division of Section~\ref{sec:dalembert} was
a pathology rather than a natural state. The compact way to put it, with apologies to Kant: symbolic
mathematics without numerical computation is impotent, and numerical computation without symbolic
mathematics is confused.

\section{Inspection of blow-up proofs by OpenAI}\label{sec:openai}

Against this stands a recent proof, produced by an AI system and formalised in Lean, that there
exists smooth, compactly supported, divergence-free initial data on $\mathbb{R}^3$ whose smooth Euler
solution has finite maximal lifespan: the velocity gradient becomes unbounded, and the
Beale--Kato--Majda~\cite{bkm} vorticity integral diverges~\cite{openaieuler}. Comparable results were being
prepared independently and by hand~\cite{buckmaster}. We take the proof to be correct, and nothing
below depends on whether it is.

It is a statement about the \emph{model}, and it is a real one: a question open for decades is
settled. Three observations place it.

\paragraph{What blows up is the gradient, and that is what turbulence does too.} The word blowup
invites the reading that the velocity becomes infinite. It does not, and the theorem does not claim
it. What is asserted is
\[
\limsup_{t\uparrow T_*}\|\grad u(t)\|_{L^\infty}=\infty ,
\qquad
\int_0^{T_*}\|\operatorname{curl}u(t)\|_{L^\infty}\,dt=\infty ,
\]
the gradient and the vorticity integral, with nothing said about the velocity becoming unbounded.

That is precisely the signature of the computed turbulent solution of
Section~\ref{sec:digital}: gradients diverging while the velocity stays bounded. The two accounts
therefore agree on the fact, that smooth solutions of \eqref{eq:mom}--\eqref{eq:slip} break down in
finite time with unbounded gradients, which was reported computationally
in~\cite{blowupHJ}. The agreement is larger than the vocabulary of blowup against turbulence
suggests, and it is worth stating before the differences.

\paragraph{But the same quantity diverges for opposite reasons.} In the turbulent flow the gradient
grows because the cascade carries energy to smaller and smaller scales and \emph{the energy leaves};
the divergence of $\grad u$ is the mechanism of dissipation, which is the drag. In the construction
the gradient grows because the energy \emph{cannot} leave --- the solution is classical and conserves
it exactly --- so a growing gradient can be accommodated only by concentrating what is there. One
divergence is energy departing, the other is energy piling up. The symptom is shared and the
physiology is opposite.

They differ also in kind of event. In the turbulent solution the unboundedness of the gradient is a
\emph{persistent condition}: the solution goes on existing, as a weak solution, dissipating, for as
long as one cares to run it. In the construction it is a \emph{terminal event}: at $T_*$ the smooth
solution ceases, and the theorem says nothing about what follows.

\paragraph{Which locates the result rather exactly.} $T_*$ is the end of the laminar regime. That is
a real thing to have established, and we would not have it dismissed: the breakdown of smooth Euler
flow in finite time was asserted on computational evidence and is now proved, for one constructed
datum. But a fluid spends essentially none of its life before $T_*$ and essentially all of it
after. The flows of Section~\ref{sec:digital} --- a wake, a separation line, a drag --- are what
happens beyond the breakdown, and they are what the equations were written to describe. The theorem
stops at the point where fluid mechanics begins.

\paragraph{It conserves energy, and therefore cannot be a fluid.} With smooth data the solution is
classical up to the blowup time, and a classical solution of \eqref{eq:mom}--\eqref{eq:inc} with
decaying data satisfies
\[
\frac{d}{dt}\int \tfrac12 |u|^2 = 0 ,
\]
by an integration by parts. The construction therefore dissipates \emph{nothing} over the whole of
its existence. The singularity is not a mechanism by which energy is lost but one by which it is
piled up: at the singular time all of it is still present, gathered into a vanishing region. A fluid
loses energy continuously from the moment it meets a body, and the loss is the drag.

\paragraph{It exhibits no individual case.} The datum is the limit of an iteration over stages and
has no closed form. Nothing can be plotted, nothing measured. Nor can either check on a \emph{result}
be applied: there is no symbolic expression to insert into
\eqref{eq:mom}--\eqref{eq:inc} and reduce, and no numerical field either, since no finite stage of
the iteration blows up --- the property belongs to the limit. What exists is the existence.

\paragraph{What is proved was already claimed.} There is a further observation, and it is not
ungenerous. The theorem asserts that \emph{some} smooth datum has finite lifespan. The position
of~\cite{blowupHJ}, stated in 2008 and resting on computation, is stronger: \emph{every} laminar Euler
flow at high Reynolds number is unstable and breaks down, with the velocity staying bounded while the
gradients do not. The existence statement follows from that as a corollary; if all smooth solutions
break down then certainly one does.

What is new is therefore the proof and not the fact. That is a genuine contribution --- a fact
asserted on computational evidence is not the same as a fact established --- and we do not diminish
it. But it should be seen for what it is: a demonstration, in a special constructed case, of a
phenomenon which is generic, which has been exhibited computationally for two decades, and which the
instability of Section~\ref{sec:digital} predicts of every flow past every body.

\paragraph{And the viscous case does not say what it is taken to say.} Two things about the
companion proof for Navier--Stokes, offered without disputing it.

The first is a matter of what was proved. Fefferman's four alternatives do not form two pairs of
negations: (A) and (B) are posed with $f\equiv0$ and quantify over \emph{all} admissible data, while
(C) and (D) ask for the existence of admissible data \emph{together with an admissible smooth
forcing} for which smoothness fails. A proof of (C) therefore leaves (A) untouched. It does not show
that a fluid tears itself into a singularity; it shows that \emph{something can be done to a fluid}
which does. Those are different statements and only the first would be about fluid flow.

The second is a matter of scale. Navier--Stokes rescales: if $v$ solves the equations at unit
viscosity with forcing $g$, then
\begin{equation}\label{eq:rescale}
u(x,t)=a\,v(bx,ct),\qquad a=\nu b,\quad c=ab,\quad f=a^2b\,g,
\end{equation}
solves them at viscosity $\nu$, and the map is onto --- one unit-viscosity construction generates the
whole family. Running it down to the viscosity of a fluid then forces a choice, and there are only
two. Hold the solution at size one, $a=1$: the forcing amplitude is $a^2b=\nu^{-1}$, so the force
required diverges as the viscosity falls and leaves any fixed admissible class. Or hold the forcing,
$a^2b=1$: then
\[
\text{amplitude } \nu^{1/3},\qquad \text{length } \nu^{2/3},\qquad \text{time } \nu^{1/3},
\]
all tending to zero --- at $\nu=10^{-6}$ the velocity reaches $10^{-2}$ on a scale of $10^{-4}$
before blowing up.

Set that beside the flow the problem is about. Turbulence at small viscosity has amplitude
\emph{fixed} and structure refining. The constructed family has amplitude \emph{vanishing} like
$\nu^{1/3}$, which is the opposite behaviour, or else a forcing diverging like $\nu^{-1}$, which is
not a fluid at all. Neither branch survives the limit and there is no third. An independent exclusion
of the mechanism under analytic forcing points the same way~\cite{civ}.

So the construction requires either a force no fluid is subject to, or a singularity too small to
constitute anything, and Horace supplies the description: \emph{parturient montes, nascetur
ridiculus mus}.

\paragraph{Forwards and backwards.} The two constructions of this paper are of the same formal kind
and opposite in direction, and it is worth saying so plainly. The singular solution fixes the
singularity wanted and solves \emph{backwards} for data and a force that will produce it, the force
being whatever the residual requires. Numerical Euler fixes the data a flow actually has and solves
\emph{forwards}, the residual becoming the bound on what may be asked of the result. Both are
automated and both exhibit something. Only the second exhibits a fluid.

\section{Three levels, and what reaches which}\label{sec:levels}

The title names three things and the point of the paper is that they are three.

\emph{A fluid flow} is a phenomenon. It has a drag which can be measured.

\emph{Euler's equations} are a model of it, parameter free, with geometry and forcing as data. They
are answerable to the phenomenon: if the drag they give is wrong, nothing can be adjusted to repair
it.

\emph{A proof about the equations} is a statement about the model. It is answerable to the model
alone.

The digital solution reaches all three levels. It is a solution of the equations, by the residual
control; it is an object that can be inspected, and inspecting such objects is where the resolution
of d'Alembert's paradox~\cite{dalembert} and the account of flight~\cite{flight} came from; and its
mean values can be set beside a measurement, where they survive.

The symbolic proof reaches the first of the three and stops. That is not a criticism of it as
mathematics, where it is an advance, and existence without an exhibited instance is ordinary and
accepted practice --- the disproof of the Mertens conjecture~\cite{mertens} established that a
counterexample exists and exhibited none. It is a statement about what the result can be offered as.
A theorem about a model becomes a statement about fluids only by way of an instance, and there is
none here.

\section{A larger split, and where it has arrived}\label{sec:split}

The divorce of Section~\ref{sec:dalembert} was not confined to fluid mechanics. It is a local
instance of a division which ran through twentieth-century mathematics and has not closed: between
symbolic, analytical, pure mathematics on one side and numerical mathematics on the other. The first
was held to be mathematics proper, concerned with existence, regularity and structure, and conducted
in symbols. The second was regarded in many departments as a service activity, a matter of machines
and approximations, and not quite the subject.

We make no case against pure mathematics, which needs none. We observe only what the division did to
\emph{this} subject. The questions the analytical side asked of Euler's and Navier--Stokes' equations
--- does a smooth solution exist, does it remain smooth, is it unique --- are questions about the
pointwise solution concept, which Section~\ref{sec:fenics} places at the stage the subject left
behind. They are hard, they are respectable, and a fluid answers none of them. Meanwhile the
quantities a fluid does determine, and which can be measured --- drag, lift, the separation line ---
were not regarded as mathematical questions at all, and the people who needed them went to a towing
tank.

The institutional expression of this is the Millennium Problem, which offers a prize for the
regularity of Navier--Stokes solutions: the apex of the analytical side, posing about fluid flow the
one question whose answer, either way, no measurement could detect.

\paragraph{And now the arrival.} That trajectory has reached a natural endpoint, and the OpenAI
proofs are it. Symbolic mathematics, equipped at last with a machine which derives faster and more
tirelessly than any person, was turned on fluid mechanics --- and what it produced, in a matter of
days, was a theorem about a datum which cannot be written down, concerning a solution which conserves
energy and therefore is not a fluid, settling a question about the pointwise concept which the
subject had already moved past. The branch has been perfected and remains disconnected. It is the
purest demonstration available of what the split costs: not that the mathematics is bad, which it is
not, but that no amount of it, done faster, reaches the thing the equations are about.

The remedy is not more symbolic power and not less. It is the collaboration of
Section~\ref{sec:fenics}: formulate symbolically, compute numerically, inspect what is computed, and
interpret symbolically what inspection shows. That loop is closed in this subject and it took two and
a half centuries to close.

\paragraph{This is not a conclusion drawn here.} The synthesis just described has been pursued as a
programme for some twenty-five years, under the name \emph{Applied Mathematics: Body and
Soul}~\cite{bodysoul,bsbooks}, and the name is the thesis: the \emph{body} of a mathematical subject
is its computational content and the \emph{soul} its analytical content, and a subject which has lost
either is crippled. The programme set out to rebuild applied mathematics and its teaching on that
union --- analysis and computation developed together from the start, rather than analysis taught as
mathematics and computation added afterwards as a skill.

The account in this paper is one instance of that programme carried through on one set of equations.
The equations are formulated symbolically; the discretisation, the residual and the adjoint are
derived symbolically; the solutions are computed; the instability is exhibited by linear analysis and
confirmed by computation; d'Alembert's paradox is resolved by inspecting what was computed; and the
result is a model which predicts drag and lift to a stated tolerance. Volume~4 of the series is the
fluid mechanics itself~\cite{book}. What the present episode supplies is a test of the thesis from
the other direction: a demonstration of what the symbolic half achieves when it is pursued, with
unprecedented power, entirely alone.

\section{The shock, and the question it raises}\label{sec:shock}

The OpenAI proofs of singularity formation for Euler and for Navier--Stokes have sent a shock wave
through the mathematical community, and the question being asked is plain enough. If a machine can
produce proofs of theorems better than people can, what is left for those whose main occupation is
proving theorems?

The question is legitimate and we do not propose to soften it. But two things about the present case
bear on it, and both follow from the argument above.

\emph{What was automated was the derivation, not the judgement.} The system worked on a problem
already posed, already named, already agreed to be worth settling, and already being attacked by
others along related lines~\cite{buckmaster}. It did not decide that finite-time blowup for Euler
was the question to ask. It did not decide that Euler's equations were the right model of a fluid,
which is a judgement about the world and not about symbols. And it did not connect its theorem to
anything: the step from a statement about a model to a statement about fluids is not a proof step,
and nothing in the output performs it. Sections~\ref{sec:openai} and~\ref{sec:levels} are an
instance of that gap rather than a complaint about it.

\emph{And this has happened before, to the discipline in this paper.} Machines outran hand
calculation in the 1950s, and numerical analysis did not end. It changed what it was about. The
object of study moved from finding a solution --- which the machine would do --- to bounding the
error in what the machine produced, to deciding what was worth computing, and to establishing when a
computed quantity is determined by the equations at all and when it is not. That is a harder subject
than the one it replaced, and it is the subject that produced the error control of
Section~\ref{sec:digital}.

We would therefore expect the answer for mathematics to have the same shape, without pretending to
know it. Deriving will be automated, and probably well. What will not be is choosing the question,
judging whether a model deserves the attention, supplying the instance which lets a theorem touch
something, and saying what a result means. On the evidence of this episode those are separate
activities from derivation, because here they were left undone.

\paragraph{And there is relevant symbolic mathematics, which is not theorem-proving.} If the proofs
discussed here concern the model and not fluids, that is no reflection on symbolic mathematics as
such. It is a reflection on one use of it. The most consequential symbolic work in physics has never
been the proving of theorems about given equations. It has been the \emph{formulation of the
equations}, from principles.

Equations \eqref{eq:mom}--\eqref{eq:slip} are themselves such a piece of work. Newton's second law is
a statement about particles; writing it for a continuum, deciding that the force of the surrounding
fluid is a pressure gradient, and adjoining incompressibility as a constraint with the pressure as
its multiplier, is symbolic mathematics of the first order --- and every step of it required deciding
what a fluid \emph{is}. The result is three lines carrying no parameters, from which everything in
this paper follows. Schr\"odinger's equation is the same kind of achievement for atoms, and the
reformulation of atomic physics as non-overlapping charge densities meeting at free interfaces, which
the present author has pursued under the name RealQM~\cite{realqm}, is an attempt at the same activity rather
than at theorem-proving: the work is in what the equations should be, not in what follows from them.

This is where symbolic mathematics does its heaviest lifting, where physical judgement enters
mathematics, and where, so far, no machine has done anything at all. A system which proves a theorem
about Euler's equations has been handed the hard part.

\section{Summary}\label{sec:summary}

Euler's equations are parameter free, with geometry and forcing as the only data, and the history of
the subject is the history of what a solution of them was taken to be. Potential flow is a pointwise
solution which exists and is unstable, so determines nothing. Leray's weak solution exists globally
and is not unique, so determines nothing either. A weak solution \emph{with its residual controlled}
is not a single flow at all: the field is not determined pointwise in space and time, but the
mean-value outputs --- drag, lift, separation --- are well determined quantities, each with a bound
the computation produces. The demand made of the solution weakens at every stage
and what can be concluded from it strengthens.

The questions symbolic analysis asks --- existence, regularity, uniqueness --- have gone unanswered
in symbolic form because turbulent solutions carry structure at every scale with nothing for
rewriting to grip, this being what turbulence is rather than a gap in technique. In digital form, by
Numerical Euler, they are answered in altered shape: solutions exist by exhibition, are not regular,
are not unique as fields, and are unique in the mean values anybody wants, with computed error bounds
and agreement with measurement.

A symbolic proof of finite-time blowup settles a long-open question about the model. Its solution
conserves energy exactly and so dissipates nothing, whereas a fluid dissipates and the dissipation is
the drag; and it exhibits no individual case, so nothing can be inserted into the equations,
inspected or measured. It is an exact statement about an object with no observables, set beside
approximate statements about observables which agree with experiment. Each has what the other lacks,
and only one of them is answerable to a fluid.

\begin{thebibliography}{99}

\bibitem{openaieuler} OpenAI, \emph{Finite time blowup for the Euler equation}, preprint, 2026,
57~pp.
\newblock \url{https://cdn.openai.com/pdf/315b36cd-ec98-4023-8342-93345194ece1/euler.pdf}

\bibitem{buckmaster} T.~Buckmaster and L.~Alp\"oge, \emph{Statement}, 2026.
\newblock \url{https://cims.nyu.edu/~tristanb/statement.pdf}

\bibitem{blowupHJ} J.~Hoffman and C.~Johnson, \emph{Blowup of incompressible Euler solutions},
BIT Numer.\ Math.\ \textbf{48} (2008) 285--307.

\bibitem{euler1755} L.~Euler, \emph{Principes g\'en\'eraux du mouvement des fluides},
M\'emoires de l'Acad\'emie des Sciences de Berlin \textbf{11} (1757) 274--315; presented 1755.
[TO VERIFY: volume and pages, and the wording of the translation quoted.]

\bibitem{bodysoul} The Body and Soul project, \url{www.bodysoulmath.org}.

\bibitem{bsbooks} K.~Eriksson, D.~Estep and C.~Johnson, \emph{Applied Mathematics: Body and Soul},
vols.~1--3, Springer, 2004; the series continuing to vol.~4 and beyond.

\bibitem{leray} J.~Leray, \emph{Sur le mouvement d'un liquide visqueux emplissant l'espace},
Acta Math.\ \textbf{63} (1934) 193--248.

\bibitem{fenics} A.~Logg, K.-A.~Mardal and G.~N.~Wells (eds.), \emph{Automated Solution of
Differential Equations by the Finite Element Method}, Lecture Notes in Computational Science and
Engineering~84, Springer, 2012.

\bibitem{book} J.~Hoffman and C.~Johnson, \emph{Computational Turbulent Incompressible Flow},
Applied Mathematics: Body and Soul, vol.~4, Springer, 2007.

\bibitem{dalembert1752} J.~le~Rond d'Alembert, \emph{Essai d'une nouvelle th\'eorie de la
r\'esistance des fluides}, Paris, 1752.

\bibitem{dalembert} J.~Hoffman and C.~Johnson, \emph{Resolution of d'Alembert's paradox},
J.\ Math.\ Fluid Mech.\ \textbf{12} (2010) 321--334.

\bibitem{flight} J.~Hoffman, J.~Jansson and C.~Johnson, \emph{New theory of flight},
J.\ Math.\ Fluid Mech.\ \textbf{18} (2016) 219--241.

\bibitem{dragcrisis} J.~Hoffman, \emph{Simulation of turbulent flow past bluff bodies on coarse
meshes using General Galerkin methods: drag crisis and turbulent Euler solutions}, Comput.\ Mech.\
\textbf{38} (2006) 390--402.

\bibitem{cs1} J.~Jansson, \emph{Adaptive Direct FEM simulation with Unicorn/FEniCS-HPC for CS1},
5th International Workshop on High Order CFD Methods; Zenodo, 2022.
\newblock \url{https://doi.org/10.5281/zenodo.6942599}

\bibitem{ifasd} J.~Jansson, \emph{Adaptive Euler: validation against NASA High Lift Prediction
Workshop~4, and preliminary tail-buffet and transonic-buffet simulations}, presented at IFASD~2026.

\bibitem{bkm} J.~T.~Beale, T.~Kato and A.~Majda, \emph{Remarks on the breakdown of smooth solutions
for the 3-D Euler equations}, Comm.\ Math.\ Phys.\ \textbf{94} (1984) 61--66.

\bibitem{kolmogorov41} A.~N.~Kolmogorov, \emph{The local structure of turbulence in incompressible
viscous fluid for very large Reynolds numbers}, Dokl.\ Akad.\ Nauk SSSR \textbf{30} (1941) 301--305.

\bibitem{taylor} G.~I.~Taylor, \emph{Statistical theory of turbulence}, Proc.\ R.\ Soc.\ Lond.\ A
\textbf{151} (1935) 421--444.

\bibitem{onsager} L.~Onsager, \emph{Statistical hydrodynamics}, Nuovo Cimento \textbf{6},
Suppl.~2 (1949) 279--287.

\bibitem{realqm} C.~Johnson, \emph{A 3D multiphase continuum model for atoms and the periodic
table}, submitted, 2026.

\bibitem{civ} P.~Constantin, M.~Ignatova and V.~Vicol, \emph{Regularity of asymptotically
axisymmetric solutions to the 3D Navier--Stokes equations with analytic forcing}, arXiv:2609.20803
(2026).

\bibitem{mertens} A.~M.~Odlyzko and H.~J.~J.~te~Riele, \emph{Disproof of the Mertens conjecture},
J.\ reine angew.\ Math.\ \textbf{357} (1985) 138--160.

\end{thebibliography}

\end{document}
